\documentclass[12pt]{article}
\usepackage[T1]{fontenc}
\usepackage[utf8]{inputenc}
\usepackage{lmodern}
\usepackage{textcomp}
\usepackage{enumitem}
\usepackage{geometry}
\geometry{margin=1in}
\begin{document}
\begin{center}
Answer Key - Science - K-12 Level Assessment 6 \
\small (Answers are ordered exactly as in the question paper.)
\end{center}

\section*{Multiple Choice}
\begin{enumerate}[label=\arabic*.]
\item \textbf{Answer:} A
\\ \textit{Options:} A) coastal ocean; B) deep ocean; C) ponds; D) lakes
\\ \textit{Note:} Most of the algae live in the coastal ocean because this environment provides the sunlight and nutrients necessary for photosynthesis, which is the process through which algae produce a significant portion of the Earth's oxygen.
\item \textbf{Answer:} A
\\ \textit{Options:} A) A car's gas mileage is increased and the engine has less power.; B) An airplane uses a more efficient engine and has higher performance.; C) A sports drink's taste is improved and has the same nutritional content.; D) A computer company upgrades the hardware and the price remains unchanged.
\\ \textit{Note:} This answer illustrates a trade-off because improving a car's gas mileage comes at the cost of reducing engine power, highlighting the principle that enhancing one feature often necessitates compromising another in design decisions.
\item \textbf{Answer:} B
\\ \textit{Options:} A) an earthquake; B) a mountain range; C) a river; D) a volcano
\\ \textit{Note:} A mountain range is formed when lithospheric plates collide and fold upon each other, resulting in the uplift of land due to tectonic forces.
\item \textbf{Answer:} A
\\ \textit{Options:} A) building a road; B) planting a tree; C) adding a freshwater source; D) creating a nature sanctuary
\\ \textit{Note:} Building a road disrupts natural habitats, fragments ecosystems, and increases pollution and human activity, which can lead to a decline in biodiversity and overall ecosystem health.
\item \textbf{Answer:} A
\\ \textit{Options:} A) converging boundaries; B) diverging boundaries; C) transform faults; D) rift zones
\\ \textit{Note:} The formation of Mount St. Helens is primarily attributed to converging boundaries where the Juan de Fuca Plate is subducting beneath the North American Plate, leading to volcanic activity.
\end{enumerate}

\section*{True / False}
\begin{enumerate}[label=\arabic*.]
\setcounter{enumi}{5}
\item \textbf{Answer:} False
\\ \textit{Options:} A) True; B) False
\\ \textit{Note:} Designing a model can help illustrate or support a hypothesis, but disproving a hypothesis typically requires empirical evidence from experiments or observations rather than just a theoretical model.
\item \textbf{Answer:} True
\\ \textit{Options:} A) True; B) False
\\ \textit{Note:} Aphids are considered first-level consumers because they obtain energy directly from plants, which are primary producers in the food web.
\item \textbf{Answer:} True
\\ \textit{Options:} A) True; B) False
\\ \textit{Note:} Light is refracted the least when passing through a tinted window because tinted windows are typically made of materials with a refractive index close to that of air, resulting in minimal bending of light compared to other optical devices that involve greater changes in refractive index.
\item \textbf{Answer:} True
\\ \textit{Options:} A) True; B) False
\\ \textit{Note:} Adding a base to a solution with a pH of 2 can increase its pH to above 7 because the base neutralizes the excess hydrogen ions present in the acidic solution, raising the pH toward neutral and potentially above it depending on the amount of base added.
\item \textbf{Answer:} False
\\ \textit{Options:} A) True; B) False
\\ \textit{Note:} The statement is false because the sun is actually an average-sized star, specifically a G-type main-sequence star, and is not considered enormous when compared to larger stars like giants and supergiants.
\end{enumerate}

\section*{Long Form Response}
\begin{enumerate}[label=\arabic*.]
\setcounter{enumi}{10}
\item \textbf{Answer:} When a fork falls off a table, the primary force acting on it is gravity. Gravity is the force that pulls objects toward the center of the Earth, causing the fork to accelerate downward as it falls. While other forces, such as air resistance, also act on the fork, they are much weaker compared to the force of gravity in this situation. As a result, the fork moves toward the floor due to the strong gravitational pull, illustrating how gravitational force influences the motion of objects.
\\ \textit{Note:} correct because it identifies gravity as the primary force acting on the fork, demonstrating how it dominates over weaker forces like air resistance to cause the fork to accelerate downward toward the floor.
\item \textbf{Answer:} Friction is crucial in bicycle ramps because it allows the tires to grip the surface, preventing slips and enabling safe acceleration and deceleration. Materials that can enhance friction include rubber, which is commonly used in tires for its high friction coefficient, and rough-textured surfaces like asphalt or concrete, which increase grip. For example, a ramp made of rubberized material would provide better traction than one made of smooth metal, as the rubber can deform slightly to create more surface contact. Additionally, using materials with added texture, such as grooves or bumps, can further enhance friction by increasing the surface area that comes into contact with the tires. Overall, selecting the right materials for bicycle ramps is essential for rider safety and performance.
\\ \textit{Note:} The answer correctly emphasizes the role of friction in ensuring safety and control on bicycle ramps, while also identifying suitable materials like rubber and rough-textured surfaces that enhance grip, thus illustrating the relationship between material choice and frictional force.
\end{enumerate}

\vfill
\noindent\textit{End of Answer Key}
\end{document}